
\section{Consequences}
    
    The observations made give rise to a multitude of interesting research questions and the observed data already contains understood behaviour like the balance of gravitational and adhesive force or the success of the half-sphere model \labelcref{eq:L_and_W_against_L0} as a tool for contact length interpretation. Crucually, the correlation between contact angles and slope parameter $k$ in the half-sphere model was confirmed for the flat-surface measurements. \\

    On the other hand, the experiments give rise to several open questions and unexpected behaviours like the crossing of contact length graphs for different cylinder diameters. To distinguish between discoveries and systematic mistakes, it is out of question that the experiment design and procedures used here can only be seen as a basis for further improvement. Therefore we would like to following criterias forwards in order to focus future research efforts on the effect of substrate geometry on droplet shapes. \\ 

    \textbf{Sample Preparation.} As the sample quality is of key interest for any high-accuracy experiments, the preparation of cylindrical samples of different materials has to be performed with high care, making sure that any form of coating does not contain any scratches or folds and preserves the cylindrical shape of the base cylinder, where it is coated on. PMMA cylinders of diameters $D \in \{3,\,4,\,5,\,7,\,10,\,15\}\,\unit{\milli\meter}$ are suggested as base cylinders for direct measurement and coating. The diameter range is chosen such that the relevant regime can be accessed using the \SI{2}{\micro\liter} - \SI{20}{\micro\liter} pipette. The coating substrates should differ substantially from PMMA in their hydrophobia. \\

    \textbf{Measurement Procedure.} While automated droplet placement is arguably the most accurate placement technique, droplet placement by hand using a high-accuracy pipette can already deliver signifiant results. The incrementation tecnique for droplet placement, as introduced in \cref{sec:Experimental_Setup}, is viewed as a method that combines reasonable accuracy and measurement at high speed. Furthermore, it allows to measure the development of a unique droplet instead of setting up a new droplet for each measurement. The tecnique is recommanded as long as each data point is averaged over sufficiently many repititions. \\

    \textbf{Setup Induced Asymmetry.} Some measurement tecniques used in this showed to affect outcomes and will therefore need improvement. Specifically, all measurements of contact widths and on flat surfaces were performed near the edge of the substrate which induce asymmetry in the droplet once the droplet touches the boundary. A better approach would be to place the droplet further away from the edge and to angle the camera downwards at a slight angle. That way, the foreground of the photo does not blure out the critical region once the camera is focussed on the contact line. Alternatively, a camera with upto 4x zoom and a higher focus depth could be used. 

\section{Conclusion}

    In this internship, the onset of water droplets on flat and cylindrical samples was investigated, comparing PMMA to PET substrate samples for cylinder diameters in the range of \SI{2}{\milli\meter} to \SI{10}{\milli\meter}. The samples were prepared, a measurement procedure was establishes and the contact lenghts and widths were measured using an optical setup. \\

    The analysis showed that contact line measurements as a function of droplet volume contribute to the same overall trajectory (\cref{fig:PMMA_L}), while contact width measurements distinguish into different graphs.\\ 
    
    Futhermore, it was found that for contact lengths at the order of the capillary length and on smaller cylinders, the droplets would rather spread along the symmetry axis of the cylinder rather than being dominated by gravity. This trend reverts for higher contact lengths and larger cylinder, such that the aspect ratio $L/W =1$ marks the overtaking of the gravitational against the adhesive force. \\
    
    The aspect ratio was analysed for different cylinder diameters and found to be highly dependent on the droplet volume such that it cannot be compared naively. By rescaling the aspect ratio against the normalized droplet volume, it could be shown that the adhesive force is not independent of the scale of the experiment, which agrees with our physical intuition. \\
    
    Lastly, we state that the contact length can be understood more easily if it is compared to the newly introduced reference length $L_0(V)$ in the half-sphere model and we suggest that its proportionality constants $k_L$ and $k_W$ could give rise to new parametrizations of material properties. \\

    All of these trends could be found for both PMMA and PET substrate, which suggests that both substrates obey similar characteristics. The comparison of proportionality constants for both paramter confirmed the expectation that PET should evolve higher contact angles than PMMA. Finally, a set of criterias for future experiment was established. 